\section{R10: forced source scale, every remote test, and a weighted local admission}

\paragraph{Scope and result.}
Retain the original weak-C2 R10 physical stationary source, every strict
source gap, $n\ge0$, the positive-infinitesimal total-scale intersection,
and every inherited exclusion. Thus $t,g,j$ are positive bounded units,
$B,h,e$ are infinitesimal, $a$ is unbounded, and $x=aB$ is infinitesimal.
The new all-tests admission proved here is the source subdomain
\[
                         n/e^2\notin\operatorname{Inf}.
\tag{R10.1}
\]
Every nonlocal independent testing chart is also proved on all of R10.
The positive-$n$, infinitesimal-$n/e^2$ local complement is handled by the following thin-layer and state-entry sections; it is not covered by this weighted pivot.
The literal $n=0$ boundary is already an inherited C1 admission, by the
exact map specified below. These statements retain all independent rates.

\paragraph{The decisive exact source scale.}
Use the pure-source variables
\[
 s=1/a,\qquad w=ae,\qquad b=a^2B=x/s,
 \quad B=s^2b,\quad x=sb,\quad e=sw,\quad r=1-sw.
\]
Initially $b$ is not assumed bounded. The gap $g=1-w+hx$ makes $w$
bounded and $1-w$ a positive unit in residue. In the exact positive
polynomials of BS2 put
\[
 Q_*=1+B(1+t+ah+2ar+at)
 +B^2\{a(h+t)+a^2[h(r+1)+t(1+r-r^2)]\},
\]
\[
 K_*=(1+s)(h+t)(1+x),\qquad \Delta_*=s\Delta.
\]
Before imposing a bound on $b$, $Q_*\sim1$ and $\Delta_*\sim1$:
write them using bounded $w$, infinitesimal $s,x,h$; their only
nonconstant monomials contain $s$ or $x$. BS1 gives exactly
\[
 \boxed{\quad 2r\Delta_*\,n/s^2=wQ_*-bK_* .\quad}
\tag{R10.2}
\]
Consequently
\[
 0<b/w\le Q_*/K_*,\qquad
 0\le n/s^2\le wQ_*/(2r\Delta_*).
\tag{R10.3}
\]
In particular $b=O(w)$ is bounded, $n=O(s^2w)$, $an$ is
infinitesimal, and $ck=\sqrt{b/t}$ is bounded for
$c=\sqrt{a/t}$, $k=\sqrt{x}$. This is an original-source bound,
including equality $n=0$, not a comparison with either testing angle.

After this bound every displayed source denominator is a unit in
$(s,b,w,h,t)$ except the explicit positive monomials already canceled.
Direct finite rational cancellation in the source formulas gives
\[
 \begin{gathered}
 m/w\sim\frac1{2(1-w)},\quad
 C\sim Ah\sim t/2,\quad D\sim\frac1{2(t-1)},\quad
 \alpha/(s^3b^2)\sim t/2,\\
 \lambda=\frac{m-a^2q}{ta^2B}\sim\frac12,\qquad
 c_F\sim\frac{t-w+tb}{2t^2}>0,\qquad
 q=nr^2+(m+\tfrac12)e^2+\alpha.
 \end{gathered}
\tag{R10.4}
\]
The last strictness follows from $t_0>1$ and $w_0<1$.
These are joint residue identities in bounded letters; they do not
set $w$ or $b$ to a positive unit. In particular they hold when either
is infinitesimal. The cancellation defining $\lambda$ leaves no $b$
or $s$ denominator; the new checker verifies this explicitly.
We also have $m$ bounded, $n,\alpha,A^{-1/2}$ infinitesimal, and
\[
 d_4\sim\frac{t}{2(t-1)},\quad b_4\sim0,\quad
 a_E\sim-\frac1{t-1},\quad x_3\sim-1.
\tag{R10.5}
\]
All first-normal arrays of the R15 source appendix are bounded:
the only newly unbounded letter $a$ occurs there through $an$ or
$aq$, both infinitesimal. The factors $a^2n$ in $c_F$ are bounded.

\paragraph{All remote tests on all of R10.}
The exact full-parent identities GR2--GR7 and WE8/WE10 apply.
Their R9 source use required $m$ infinitesimal, but their algebra does
not. Under R10.4 the entire nonnegative $m$ packet is retained.
For GR2 it adds a positive rank-one form to the same strict six-root
residue GR3; all its coefficients are bounded. For the chart
$u$ infinitesimal, $v$ a bounded unit, keep $S=\sqrt m W$.
The complete $(S,X)$ residue is now
\[
 (v_0S-\sqrt{m_0}X)^2+\frac{1+v_0^2}{2}X^2,
 \qquad \det=\frac{v_0^2(1+v_0^2)}2>0.
\tag{R10.6}
\]
The $(uW,Z)$ block remains GR5 and the normalized $Y$ pivot remains
strict; the extra $uvY,uZ$ terms in the $m$ packet are infinitesimal.
For unbounded $v$ its reciprocal identity has the same diagonal
residue GR6 because the potential $\sqrt m X$ term carries $1/v$.
On the remaining remote charts retain $mR_m^2$ as a full positive
square and apply the unchanged strict residue to the remainder.
The original homogeneous identities include either literal infinity,
and the GR2/WE8 axes retain $n=0$ without a division. The original
sign gauges and real/imaginary split are unchanged. Hence every test
outside both angles infinitesimal is covered for the entire R10 source.

\paragraph{The exact local full-parent identity.}
Use the R15 generic contact balances RC1, RC6 and exact normal
identities RC7--RC9. These are rational identities of the original
source and never require bounded $a$. Set $U=u^2,V=v^2$, $y=tB$,
\[
 E=tuW-vZ,\quad N_y=(Y-huW)/\sqrt h,\quad \xi=X+ruZ,
\]
\[
 F_v=(V-y)Z,\quad R_D=v(U-x)Z,\quad J=(uv-x/c)Z,
\]
\[
 Q=\sqrt q(avW-uZ-a\xi),\qquad
 N=\sqrt n\{r[(a/t)V-U]Z+u\xi\},
\]
\[
 R=\sqrt\lambda\,tk(cv-u)W,\qquad
 T=\sqrt{2\lambda/c}\,k\sqrt{uv}(tcW-Z),\quad
 \mu=m-a^2q=ta\lambda x.
\]
At positive nonzero angles these are exactly the R15 coupled roots.
They have no testing-ratio inverse. Let $\mathcal H_0$ denote the
hard quadratic after subtracting $Q^2+N^2$ and the exact rank packet
$\mu\{vE/(tu)-\xi\}^2$ from the original parent. Independent
full-parent differentiation in the new checker verifies that every
entry of $\mathcal H_0$ has no $a$ dependence and no testing inverse.
Its R10 local residue is
\[
 \operatorname{diag}\left(\frac1{2(t_0-1)},\frac{t_0}2,\frac12\right)
 \quad\hbox{on }(E,N_y,\xi).
\tag{R10.7}
\]
This permits bounded $m$; every occurrence of $m$ in the remaining
hard entries carries a positive testing degree or $h$.

The principal angular--$\xi$ coupling is resolved as a whole.
Writing $R_c=kv(tcW-Z)$, direct multiplication gives
\[
 -2t\lambda ck R_c\xi+2t\lambda c(R_D-uJ)\xi
 =-2t\lambda x Q\xi/\sqrt q-2\mu\xi^2.
\tag{R10.8}
\]
Thus neither $cu$ nor $cR_D$ is left as a Gram coefficient. Combining
RC2, RC7--RC10, the full rank packet and R10.8 gives the entire
original $F/4$ exactly as
\[
\begin{split}
 &Q^2+N^2+T^2+R^2+\lambda J^2+d_4R_D^2
 +(c_F+b_4x/t^2)F_v^2+(b_4/t^2)vR_DF_v\\
 &\quad+\mathcal H_0-\lambda xE^2-\mu\xi^2
                  -2(t\lambda x/\sqrt q)Q\xi\\
 &\quad+E\{(a_E+c_EV)R_D+(b_E+c_Ex)vF_v\}\\
 &\quad+\sqrt hN_y\{(a_Y+c_YV)R_D+(b_Y+c_Yx)vF_v\}\\
 &\quad+\xi\{x_5uvR_D+u(x_3+x_5x)F_v\}.
\end{split}
\tag{R10.9}
\]
All coefficients refer to the exact source arrays in the R15 appendix.
This identity retains every scalar, first-normal and hard-normal term.

\paragraph{A strict weighted pivot when $n/e^2$ is noninfinitesimal.}
Put $\rho=n/e^2$, $H=m+1/2$, $\chi=m/w$, and
$\overline q=a^2q$. Then
\[
 \overline q=w^2(\rho r^2+H)+a^2\alpha,
 \qquad \chi-H\in\operatorname{Inf}.
\tag{R10.10}
\]
The $Q,\xi$ source block in R10.9 has Schur pivot
\[
 S=\frac12-\mu-\frac{(t\lambda x)^2}{q}
   =H-\frac{m^2}{\overline q}.
\tag{R10.11}
\]
The equality uses $a^2q+\mu=m$ exactly. Since $\alpha\ge0$,
\[
 S\ge H-\frac{\chi^2}{\rho r^2+H}
 =\frac{H\rho r^2}{\rho r^2+H}
    +\frac{H^2-\chi^2}{\rho r^2+H}.
\tag{R10.12}
\]
Here $H$ is a positive bounded unit. Under R10.1 the first summand
has positive residue and the second is infinitesimal, also when
$\rho$ is unbounded. Hence $S$ is a positive bounded unit.
Moreover $b/w$ is bounded and $\overline q\ge Hw^2$, so
$x/\sqrt q=b/\sqrt{\overline q}$ is bounded. All entries of this
strict source block are therefore bounded.

In the rest of R10.9 the only nonzero off-diagonal residue is
$a_E$ between $E$ and $R_D$, giving
\[
 \frac1{2(t_0-1)}
 \begin{pmatrix}1&-1\\-1&t_0\end{pmatrix}\succ0.
\]
The $N_y,F_v,J,R$ pivots have positive residues
$t_0/2,(t_0-w_0+t_0b_0)/(2t_0^2),1/2,1$.
Every other entry is infinitesimal by R10.4--R10.5 and the explicit
positive testing degrees in R10.9. Retain $N^2,T^2$ as exact positive
squares. The bounded-Gram lemma supplies finite tensor SOS for the
entire parent on R10.1. No source/testing rate comparison occurs.
All roots homogenize without $Z$ or testing-ratio inverses, so literal
axes and $Z=0$ are included directly. The original signs, complex-state
split and the remote theorem complete the all-tests admission.

\paragraph{Literal $n=0$ and the exact remaining complement.}
At $n=0$, OP4 is a physical C1 identification with
\[
 x_1=x,\quad\eta_1=t,\quad\zeta_1=s,\quad r_1=h,
 \quad\theta_1=e,\quad e_{1,\mathrm{scale}}=\sqrt{tb}.
\]
Thus $k_1,\zeta_1,r_1$ are infinitesimal, $\eta_1$ a unit, and
$e_{1,\mathrm{scale}}$ bounded. Its infinitesimal type is covered by
\texttt{baseline/baseline/baseline/advances/c1\_scales/small\_e\_all\_tests.tex};
its positive-unit type is covered by the neighboring
\texttt{centered\_a\_bounded\_delta.tex}. Both retain $\beta=0$.
This is C1.2/C1.3 routing, not the unbounded-$\eta$ C1.1 theorem.
For $n>0$ OP5 gives formal $\beta_1<0$, so this physical C1 theorem
cannot be extended by continuity or imported into the neighborhood.
The complement left for the companion thin-layer proof is exactly
\[
       0<n/e^2\in\operatorname{Inf},\qquad u,v\in\operatorname{Inf},
\]
with every original physical and total-scale intersection retained.
In this complement R10.11 can have zero residue: at a unit $w$ and
$b_0=w_0/t_0$, its residue is zero. Thus R10.9 is a useful full-parent
identity there, but its stated strict-Gram argument does not prove that
complement. The companion thin-layer proof and direct state join complete this complement; the final join states the whole-R10 theorem.
